\chapter{Finite-Depth Quantum Stability and Ultraviolet Growth}
\label{ch:radiative}

\section{Why finite depth is not optional}

The classical fiber operator is well defined on suitable infinite-depth domains.  A quantum effective potential asks a different question.  Every level contributes vacuum fluctuations, and the multiplicity grows with depth.  Normal ordering a produced-particle density does not remove the vacuum terms in the effective action.

The honest starting point is therefore finite $H$ with a declared physical or effective cutoff.  Let
\begin{equation}
M_h^2(\phi)=m_b^2+F(\phi)\lambda_h,
\qquad
\lambda_h=\Lambda^2p^{2\sigma h}.
\label{eq:rad-level-mass}
\end{equation}

\section{One-loop effective potential}

For canonical real bosonic detail fields on a constant background with $M_h^2(\phi)>0$, the flat-space, local-potential contribution in a standard subtraction scheme is
\begin{equation}
\boxed{
\Delta U_H(\phi)
=\frac1{64\pi^2}
\sum_{h=1}^{H}d_hM_h^4(\phi)
\left[
\log\frac{M_h^2(\phi)}{\bar\mu^2}-\frac32
\right].}
\label{eq:rad-CW}
\end{equation}
At the locked origin, $F'(0)=0$ and $F''(0)=2\gamma/v^2$.  The induced average curvature is
\begin{equation}
\boxed{
\Delta m_\phi^2
=\frac\gamma{16\pi^2v^2}
\sum_{h=1}^{H}d_h\lambda_hM_{h0}^2
\left[
\log\frac{M_{h0}^2}{\bar\mu^2}-1
\right],}
\label{eq:rad-mass-correction}
\end{equation}
where $M_{h0}^2=m_b^2+\lambda_h$.

Equation \eqref{eq:rad-mass-correction} does not by itself prove that locking fails.  Counterterms can be chosen at finite depth.  It proves that the locking curvature is radiatively sensitive to the tower unless a symmetry or ultraviolet completion controls the sum.  On a curved background, additional $M_h^2R$ and curvature-squared counterterms are required.  If a level mass becomes negative, the logarithm develops an imaginary part and the expression is not a real stability potential.

\section{Infinite-depth growth theorem}

\begin{theorem}[Bosonic tower ultraviolet growth]
Fix $\phi$ with $F(\phi)>0$.  For $\sigma>0$ and large depth,
\begin{equation}
\sum_{h=1}^{H}d_hM_h^4
\asymp
F(\phi)^2\Lambda^4p^{(1+4\sigma)H}.
\label{eq:rad-geometric-growth}
\end{equation}
Writing
\begin{equation}
M_{\mathrm{UV}}=\Lambda p^{\sigma H},
\end{equation}
the leading power is
\begin{equation}
\boxed{
\Delta U_H
\sim
\frac{M_{\mathrm{UV}}^{4+1/\sigma}}
{\Lambda^{1/\sigma}}}
\label{eq:rad-UV-power}
\end{equation}
up to logarithms and scheme-dependent coefficients.
\end{theorem}

\begin{proof}
At high level, $d_h\sim p^h$ and $M_h^4\sim F(\phi)^2\Lambda^4p^{4\sigma h}$.  Their product is a geometric sequence with ratio $p^{1+4\sigma}$.  The last term controls the finite sum.  Substituting $p^H=(M_{\mathrm{UV}}/\Lambda)^{1/\sigma}$ gives \eqref{eq:rad-UV-power}, with the fixed factor $F(\phi)^2$ suppressed there.  If $F=0$, the base mass instead gives growth proportional to $p^H$.  If $F<0$, sufficiently high levels are tachyonic.
\end{proof}

\begin{figure}[htbp]
\centering
\includegraphics[width=0.82\linewidth]{radiative_uv_growth.pdf}
\caption{The multiplicity-weighted fourth mass moment at finite depth.  The fitted slopes agree with $4+1/\sigma$.  The lines have different horizontal ranges because the same depth produces different ultraviolet masses.}
\label{fig:radiative-growth}
\end{figure}

\begin{negativebox}[Ultraviolet boundary]
The ordinary bosonic infinite tower does not have a finite one-loop potential without a physical cutoff, cancellations, or a different spectral ultraviolet prescription.  Classical Hilbert-space convergence does not supply that cancellation.
\end{negativebox}

\section{Radiative closure of the one-coupling pattern}

At tree level, the level interaction coefficients obey
\begin{equation}
g_h=\gamma\frac{\lambda_h}{v^2}.
\end{equation}
A generic finite multiscalar quantum theory produces counterterms of the forms
\begin{equation}
\phi^4,
\qquad
\chi_h^4,
\qquad
\chi_h^2\chi_j^2,
\qquad
\phi^2\chi_h^2,
\end{equation}
with independent finite parts unless a symmetry relates them.  Level-orthogonal invariance constrains indices within one level.  Branch homogeneity relates levels only if the regulator and renormalization prescription preserve the branching maps.  A finite upper boundary introduces additional endpoint terms.

The single $\gamma$ statement applies only to the order-$2\sigma$ spectral channel of the classical ansatz.  The allowed order-zero coefficient $\kappa$ was set to zero as a model restriction, and generic loop corrections need not preserve that restriction.  There is no renormalization theorem yet.  A completed model needs one of the following:
\begin{enumerate}[leftmargin=2.2em]
\item a physical finite-depth endpoint
\item paired degrees of freedom that cancel the leading vacuum moments
\item an ultraviolet fixed point with a proved invariant counterterm family
\item a modulus or gauge principle that turns the coefficient pattern into a protected geometric quantity
\end{enumerate}

\section{A dynamical refinement-scale modulus}

One candidate completion promotes the multiplier of $L_H$ to a field.  Introduce a canonical scalar $s$ and define
\begin{equation}
\begin{aligned}
S_H^{\mathrm{mod}}=\int\sqrt{-g}\,\dd^4x\Bigg[
&\frac{M_{\mathrm{Pl}}^2}{2}R
-\frac12(\nabla\phi)^2-U(\phi)
-\frac12(\nabla s)^2-W(s,\phi)\\
&-\frac12\ip{\nabla\chi}{\nabla\chi}_H
-\frac12\ip{\chi}{\left(m_b^2+\ee^{2s/f_s}L_H\right)\chi}_H
\Bigg].
\end{aligned}
\label{eq:rad-modulus-action}
\end{equation}
Choose
\begin{equation}
W(s,\phi)
=\frac{M_s^2}{2}
\left[
s-\frac{f_s}{2}\log F(\phi)
\right]^2.
\label{eq:rad-modulus-potential}
\end{equation}

\begin{proposition}[Heavy-modulus matching]
If
\begin{equation}
M_s^2\gg H^2,
\qquad
M_s^2\gg A_\star m_{\mathrm{peak}},
\qquad F>0,
\qquad
|\Box s_0|\ll f_sM_s^2,
\qquad
F Y_H\ll f_s^2M_s^2,
\label{eq:rad-heavy-conditions}
\end{equation}
where $s_0(\phi)=\tfrac{f_s}{2}\log F(\phi)$.  Then the modulus equation has the adiabatic solution
\begin{equation}
s=s_0
+\frac{\Box s_0}{M_s^2}
-\frac{FY_H}{f_sM_s^2}
+\order(M_s^{-4}),
\label{eq:rad-s-solution}
\end{equation}
and hence
\begin{equation}
\ee^{2s/f_s}
=F(\phi)
\left[
1+\order\left(\frac{\Box s_0}{f_sM_s^2},
\frac{FY_H}{f_s^2M_s^2}\right)
\right].
\label{eq:rad-modulus-match}
\end{equation}
\end{proposition}

\begin{proof}
Expand the modulus equation around the moving minimum of \eqref{eq:rad-modulus-potential}.  To leading order it is $\Box s_0-M_s^2(s-s_0)-FY_H/f_s=0$, which gives \eqref{eq:rad-s-solution}.  Expanding the exponential gives \eqref{eq:rad-modulus-match}.
\end{proof}

\begin{statusbox}[Model status]
The proposition matches the detail mass multiplier, not the complete original scalar EFT.  Substituting the moving minimum also produces the unsuppressed operator
\begin{equation}
-\frac12(\nabla s_0)^2
=-\frac{f_s^2}{8}(\nabla\log F)^2.
\end{equation}
It must be retained or shown negligible.  The chosen potential \eqref{eq:rad-modulus-potential} remains a model input.  It does not derive $\gamma$, and it does not cancel the tower vacuum energy by itself.
\end{statusbox}

\section{Quantum tasks that remain}

The next quantum block must supply a regulator compatible with diffeomorphism and internal symmetries, a finite counterterm basis, a measure contribution in the BV pushforward, and a cutoff-stable prediction.  It must also decide whether $H$ is physical, emergent, or removed through a cancellation.  Until then, infinite-depth statements should be restricted to the classical operator domains for which they were proved.
